\begin{proofbox}
	\textbf{\textcolor{CProof}{思路提示}}

	利用椭圆参数方程将问题转化为三角函数值域问题，再通过辅助角公式求出
	\[
		M=Aa\cos\theta+Bb\sin\theta
	\]
	的取值范围，最后分情况讨论绝对值即可。

	\medskip

	\textbf{\textcolor{CProof}{正式推导}}

	\begin{description}[style=nextline, leftmargin=2.8em, labelsep=0.5em, itemsep=0.55em]
		\item[\textbf{步骤一（参数化距离）：}]
		      设椭圆 \(E\) 上点为
		      \[
			      (a\cos\theta,\,b\sin\theta),
			      \qquad \theta\in[0,2\pi).
		      \]
		      点到直线 \(l:Ax+By+C=0\) 的距离为
		      \[
			      d(\theta)
			      =
			      \frac{|Aa\cos\theta+Bb\sin\theta+C|}{\sqrt{A^2+B^2}}.
		      \]

		      \par\textbf{理由：}
		      椭圆的参数方程形式统一，代入点到直线的距离公式即可得到关于 \(\theta\) 的函数。

		\item[\textbf{步骤二（求分子内部的值域）：}]
		      令
		      \[
			      M(\theta)=Aa\cos\theta+Bb\sin\theta.
		      \]
		      利用辅助角公式：
		      \[
			      M(\theta)
			      =
			      \sqrt{a^2A^2+b^2B^2}\,\sin(\theta+\varphi)
			      =
			      R\sin(\theta+\varphi),
		      \]
		      其中
		      \[
			      R=\sqrt{a^2A^2+b^2B^2},
		      \]
		      并可取
		      \[
			      \cos\varphi=\frac{Bb}{R},
			      \qquad
			      \sin\varphi=\frac{Aa}{R}.
		      \]
		      所以
		      \[
			      -R\le M(\theta)\le R.
		      \]

		      \par\textbf{理由：}
		      辅助角公式将两个三角函数合并为一个正弦函数，其值域由系数平方和开方决定。

		      进一步地，
		      \[
			      M(\theta)=R
		      \]
		      在
		      \[
			      \cos\theta=\frac{aA}{R},
			      \qquad
			      \sin\theta=\frac{bB}{R}
		      \]
		      时取得，对应椭圆上的点为
		      \[
			      P_+=\left(\frac{a^2A}{R},\frac{b^2B}{R}\right).
		      \]
		      同理，
		      \[
			      M(\theta)=-R
		      \]
		      在
		      \[
			      P_-=-P_+
		      \]
		      处取得。

		\item[\textbf{步骤三（讨论 \(|M+C|\) 的最值）：}]
		      记
		      \[
			      N(\theta)=|M(\theta)+C|.
		      \]

		      \begin{itemize}[leftmargin=*]
			      \item 当 \(|C|>R\) 时，区间 \([C-R,C+R]\) 不含 \(0\)，因此 \(M(\theta)+C\) 恒与 \(C\) 同号，故
			            \[
				            N_{\min}=|C|-R,
				            \qquad
				            N_{\max}=|C|+R.
			            \]

			      \item 当 \(|C|\le R\) 时，存在 \(\theta\) 使
			            \[
				            M(\theta)=-C,
			            \]
			            所以
			            \[
				            N_{\min}=0.
			            \]
			            同时端点到 \(0\) 的最大距离为
			            \[
				            N_{\max}=|C|+R.
			            \]
		      \end{itemize}

		      \par\textbf{理由：}
		      因为 \(M\) 的范围为 \([-R,R]\)，所以 \(M+C\) 的范围为
		      \[
			      [C-R,C+R].
		      \]
		      对这个区间整体取绝对值，最小值取决于区间是否包含 \(0\)，最大值一定等于端点到 \(0\) 的较大距离，即 \(|C|+R\)。

		\item[\textbf{步骤四（回归距离公式）：}]
		      将 \(N(\theta)\) 的最值除以 \(\sqrt{A^2+B^2}\)，即得
		      \[
			      d_{\min}
			      =
			      \max\left\{0,\frac{|C|-R}{\sqrt{A^2+B^2}}\right\},
			      \qquad
			      d_{\max}
			      =
			      \frac{|C|+R}{\sqrt{A^2+B^2}}.
		      \]

		      \par\textbf{理由：}
		      距离公式的分母 \(\sqrt{A^2+B^2}\) 与 \(\theta\) 无关，因此距离的最值直接对应分子最值除以常数。
	\end{description}

	\medskip

	\textbf{\textcolor{CProof}{方法二（独立证明：切线法向量与平行线距离）}}

	\textbf{思路提示：}固定已知直线的方向，在椭圆两侧找到与它平行的两条切线；
	它们是沿该法向方向能够到达的两个边界。先求切点，再写切线方程，最后比较平行线间的距离。

	\begin{description}[style=nextline, leftmargin=2.8em, labelsep=0.5em, itemsep=0.55em]
		\item[\textbf{步骤一（让切线的法向量与已知直线同向或反向）：}]
		      设切点为 \(P(x_0,y_0)\)，椭圆在此点的切线方程为
		      \[
			      \frac{x_0x}{a^2}+\frac{y_0y}{b^2}=1.
		      \]
		      切线的法向量为
		      \[
			      \vec n_P=\left(\frac{x_0}{a^2},\frac{y_0}{b^2}\right),
		      \]
		      已知直线 \(l:Ax+By+C=0\) 的法向量为 \((A,B)\)。
		      两条直线平行，当且仅当其法向量成比例，故设
		      \[
			      \left(\frac{x_0}{a^2},\frac{y_0}{b^2}\right)
			      =\lambda(A,B),\qquad \lambda\ne 0.
		      \]
		      从而
		      \[
			      x_0=\lambda a^2A,\qquad y_0=\lambda b^2B.
		      \]

		      \par\textbf{理由：}一般切点处的切线方向随切点变化。
		      我们只寻找其中与原直线平行的切线，所以直接约束二者法向量平行。
		      法向量可以反向，因此必须保留比例系数的正、负两种可能。

		\item[\textbf{步骤二（联立椭圆方程，求出两个切点）：}]
		      因为切点在椭圆上，代入
		      \[
			      \frac{x_0^2}{a^2}+\frac{y_0^2}{b^2}=1
		      \]
		      得
		      \[
			      \lambda^2a^2A^2+\lambda^2b^2B^2=1,
		      \]
		      即
		      \[
			      \lambda^2R^2=1,\qquad
			      \lambda=\pm\frac1R,
		      \]
		      其中 \(R=\sqrt{a^2A^2+b^2B^2}>0\)。
		      因此
		      \[
			      \boxed{P_\pm=\pm\left(\frac{a^2A}{R},\frac{b^2B}{R}\right)}.
		      \]

		      \par\textbf{理由：}同一个法向方向对应椭圆相对两侧的两个支撑点；
		      它们关于椭圆中心成中心对称。

		\item[\textbf{步骤三（将两个切点代回切线方程）：}]
		      在 \(P_+\) 处，切线方程为
		      \[
			      \frac{(a^2A/R)x}{a^2}+\frac{(b^2B/R)y}{b^2}=1,
		      \]
		      化简得
		      \[
			      l_+:Ax+By-R=0.
		      \]
		      同理，在 \(P_-\) 处得到
		      \[
			      l_-:Ax+By+R=0.
		      \]
		      所以椭圆夹在这两条平行的支撑切线之间，即
		      \[
			      -R\le Ax+By\le R.
		      \]

		      \par\textbf{理由：}切线是椭圆沿给定法向方向的边界。
		      两条切线之间的宽度为 \(2R/\sqrt{A^2+B^2}\)，
		      且椭圆的整个边界都位于这条带状区域内。

		\item[\textbf{步骤四（平行直线间距离，分类求最值）：}]
		      将原直线 \(l:Ax+By+C=0\) 分别与 \(l_+,l_-\) 比较，得
		      \[
			      D_+=\frac{|C+R|}{\sqrt{A^2+B^2}},\qquad
			      D_-=\frac{|C-R|}{\sqrt{A^2+B^2}}.
		      \]
		      最远点必在两支撑点之一，因此
		      \[
			      d_{\max}=\max\{D_+,D_-\}
			      =\frac{|C|+R}{\sqrt{A^2+B^2}}.
		      \]

		      若 \(|C|>R\)，原直线在两条支撑切线之外，故
		      \[
			      d_{\min}=\min\{D_+,D_-\}
			      =\frac{|C|-R}{\sqrt{A^2+B^2}}.
		      \]
		      若 \(|C|=R\)，原直线就是其中一条支撑切线，故 \(d_{\min}=0\)。
		      若 \(|C|<R\)，原直线穿过两支撑切线夹成的带状区域，
		      与椭圆相交于两点，故同样有 \(d_{\min}=0\)。

		      \par\textbf{理由：}两条支撑切线可以直接确定\emph{最远}位置；
		      但当原直线穿过椭圆时，\emph{最近}点是交点，
		      不能把原直线到两条支撑切线的较小距离误认为最小点线距离。
	\end{description}

	\medskip

	\textbf{\textcolor{CProof}{结论回扣}}

	\textbf{由以上推导可见，}
	\[
		R=\sqrt{a^2A^2+b^2B^2}
	\]
	\textbf{是核心辅助量。它控制了 \(Ax+By\) 在椭圆上的取值范围，也就完全控制了椭圆上点到直线距离的最值。}

\end{proofbox}